\scanpage{245}
\mspage{195}

muni d'un morphisme

\[
i':\mathrm{id}_{\mathrm{Cat}_{/S}}\longrightarrow C'
\]

qui est dans $W$ argument par argument, et tel que $C'$ se factorise par $\underline{\mathrm{Fib}}_W S$.

On peut comparer $C$ et $C'$ moyennant $C\circ C'$, car

\diagram{\begin{tikzcd}[ampersand replacement=\&]
\&C\circ C'\&\\
C\unsure{(}{}=C\circ\mathrm{id}\arrow[ur,very thick,"{C*i'}"]\arrow[ur,phantom,"{\deleted{\unsure{\swarrow}}}"{pos=.18,swap}]\&\&C'=\mathrm{id}\circ C'\arrow[ul,"{i*C'}"']
\end{tikzcd}}

i.e. on a, pour tt $X$ dans $\mathrm{Cat}_{/S}$

\margin{\unsure{''}}

\diagram{\begin{tikzcd}[ampersand replacement=\&]
\&C(C'(X))\&\\
C(X)\arrow[ur,"{C(i'_X)}"]\&\&C'(X)\arrow[ul,"{i_{C'(X)}}"'],
\end{tikzcd}}

ces deux flèches étant dans $W$ (celle de gauche parce que $C(W)\subset W$). Ce sont des flèches dans $\underline{\mathrm{Fib}}_W S$, et si on \deleted{\unsure{admet}} suppose que $W$ satisfait $\unsure{\mathrm{W}(4|\mathrm{WL}4|\mathrm{W}4}\,\mathrm b)$, elles induisent des $W$-équivalences en chaque fibre. Cela montre donc déjà
% archive sheet number: 245
